
\documentclass[12pt,a4paper]{article}

\usepackage[T1]{fontenc}
\usepackage{lmodern}
\usepackage{xcolor}

\title{Colors and Boxes in LaTeX}
\author{Vikram Mehra}
\date{\today}

% Define custom colors
\definecolor{myblue}{RGB}{0,70,150}
\definecolor{mygreen}{RGB}{0,130,70}
\definecolor{lightyellow}{RGB}{255,245,180}

\begin{document}

\maketitle

\section{Introduction}

LaTeX supports colored text, backgrounds, borders,
and customized boxes using the xcolor package.

\section{Changing Text Color}

The textcolor command changes the color of selected text.

\textcolor{red}{This is red text.}

\textcolor{blue}{This is blue text.}

\textcolor{mygreen}{This is custom green text.}

\section{Defining Custom Colors}

The definecolor command creates a new color.

\textcolor{myblue}{This text uses a custom blue color.}

\textcolor{mygreen}{This text uses a custom green color.}

\section{Creating a Colored Background}

The colorbox command places a colored background
behind text.

\colorbox{yellow}{Text with a yellow background.}

\bigskip

\colorbox{lightyellow}{Text with a custom background.}

\section{Creating a Colored Border}

The fcolorbox command creates a box with a colored
border and background.

\bigskip

\fcolorbox{blue}{yellow}{LaTeX Colored Box}

\bigskip

\fcolorbox{myblue}{lightyellow}{
    \parbox{0.7\textwidth}{
        This is a customized box containing multiple
        lines of text. It demonstrates how LaTeX can
        combine colored borders and backgrounds.
    }
}

\section{Changing the Text Color of a Section}

{\color{myblue}
This complete paragraph is displayed in custom blue.
The color declaration remains active within its braces.
}

\section{Conclusion}

The xcolor package provides useful commands for
customizing documents with text colors, background
colors, and colored boxes.

\end{document}
